\scanpage{185}
% archive foliation: 184
mais il est naturel de mettre sur $\pi_0(L)$ la topologie image inverse de
$K$ topologie \inserted{des $\pi_0(L)$} [des partitions ouvertes de $L$].
$F$ est dit \uline{connexe} si $F\simeq\mathrm{id}_{\mathrm{Ens}}$,
i.e. $F$ est $T$-représentable par l'espace ponctuel
[i.e. par un espace connexe].

Considérons deux \deleted{foncteurs} flèches
\diagram{$F^*:$\quad\begin{tikzcd}
F^0\arrow[r,shift left=.5ex]\arrow[r,shift right=.5ex]&F^1
\end{tikzcd}
\qquad (F^0,F^1\in\underline{\operatorname{Hom}}(\mathrm{Ens},\mathrm{Ens}))}
et le \uline{foncteur noyau}
\[
\Pi^0(F^*)=F=\operatorname{Ker}(F^0\rightrightarrows F^1)
\]
Si $F^0,F^1$ sont \deleted{\unsure{strictement pro}}\inserted{$T$-}représentables,
alors les deux flèches proviennent de flèches
\diagram{\begin{tikzcd}
K_0&K_1\arrow[l,shift left=.5ex]\arrow[l,shift right=.5ex]
\end{tikzcd}}
et $F$ sera également $T$-représentable
[N.B. Toute $\varprojlim$ de foncteurs $T$-représentables est $T$-repr,
en prenant une limite inductive sur les espaces dont ils proviennent,
\deleted{\unsure{\xi}} et le $\pi_0$ associé].
Un couple de $F^{\unsure{\alpha|*}}$ comme dessus est dit
\uline{connexe} si son noyau est connexe.

Considérons un ``foncteur \unsure{cosimplicial}'',
\inserted{\unsure{objet simplicial de}}
de $\underline{\operatorname{Hom}}(\mathrm{Ens},\mathrm{Ens})$
\diagram{\begin{tikzcd}
F^0\arrow[r,shift left=.5ex]\arrow[r,shift right=.5ex]&
F^1\arrow[r,shift left=1ex]\arrow[r]\arrow[r,shift right=1ex]&F^2
\end{tikzcd}}
On lui associe foncteurs $\Pi^0(F^*;E)$,
\deleted{\unsure{G}} et $\Pi^1(F^*;G)$,
où $E$ ens variable, $G$ groupe variable.
Si $F^*$ est \uline{coponctué} i.e. muni de
$F^*\to\mathrm{id}_{\mathrm{Ens}}^*$, on définit des $=$
$\Pi^0(F^*,\unsure{\xi|\zeta};E)$ et
$\Pi^1(F^*,\unsure{\xi|\zeta};G)$, $\unsure{\xi|\zeta}$
